% --- Archon physics formalization source begin ---
% archon:physics
% archon:covers PhyXMiniProblems/problem_phyx_mini_0498.lean
% archon:source-report reports/phyx_mini/problem_phyx_mini_0498.source.json

\chapter{Physics problem phyx\_mini\_0498}
\label{ch:PhyXMiniProblems_problem_phyx_mini_0498}

\paragraph{Problem source.}
\#\# Question

How many photoelectrons are ejected per second in the ex-periment represented by the figure?

\#\# Answer choices

- A: A: \textbackslash{}( 6.25 \textbackslash{}times 10\textasciicircum{}\{12\} \textbackslash{})

- B: B: \textbackslash{}( 6.25 \textbackslash{}times 10\textasciicircum{}\{14\} \textbackslash{})

- C: C: \textbackslash{}( 6.25 \textbackslash{}times 10\textasciicircum{}\{13\} \textbackslash{})

- D: D: \textbackslash{}( 6.25 \textbackslash{}times 10\textasciicircum{}\{16\} \textbackslash{})

\#\# Figure caption (auxiliary; use the image as primary evidence)

This is a line graph depicting the relationship between stopping potential (\textbackslash{}(V\_\{\textbackslash{}text\{stop\}\}\textbackslash{})) and frequency (\textbackslash{}(f\textbackslash{})).

- **Axes**:
  - The vertical axis (y-axis) represents the stopping potential, labeled as \textbackslash{}(V\_\{\textbackslash{}text\{stop\}\} \textbackslash{}, (\textbackslash{}text\{V\})\textbackslash{}), with a range from 0 to 9 volts.
  - The horizontal axis (x-axis) represents the frequency, labeled as \textbackslash{}(f \textbackslash{}, (\textbackslash{}times 10\textasciicircum{}\{15\} \textbackslash{}, \textbackslash{}text\{Hz\})\textbackslash{}), with a range from 0 to 3.

- **Line**:
  - A green line extends diagonally from the origin (0,0) to the point (3,9), indicating a linear relationship.

- **Grid**:
  - The graph includes a grid for better visualization of data points, with equal intervals marked along both axes.

The graph likely represents the relationship between the frequency of incident light and the resulting stopping potential in a photoelectric effect experiment, showing a direct proportionality.

\#\# Dataset metadata

Quantum Phenomena

\paragraph{Recorded answer/context.}
C: C: \textbackslash{}( 6.25 \textbackslash{}times 10\textasciicircum{}\{13\} \textbackslash{})

\paragraph{Figure/image path.}
/root/proposal\_for\_physic/science-mango/phyx\_mini\_run/phyx\_data/test\_image/498.png

\paragraph{Formalization target.}
create a compiling Lean file with sorry bodies at `PhyXMiniProblems/problem\_phyx\_mini\_0498.lean`. The Lean declarations must preserve the physical quantities, dimensions or dimensional roles, figure labels, governing-law hypotheses, and final relation expressed by this problem.
Use Mathlib/Physlib names found through LeanExplore where available. If a physics API is missing, introduce faithful local abstractions rather than scalar placeholder aliases.

\begin{theorem}[Physics formalization target]
\label{thm:physics:phyx_mini_0498:target}
\lean{PhyXMiniProblems.ProblemPhyXMini0498.problem_phyx_mini_0498}
\uses{def:physics:phyx-mini-0498:phyxminiproblems-problemphyxmini0498-photoelectronratepersecond, def:physics:phyx-mini-0498:phyxminiproblems-problemphyxmini0498-photoelectricexperiment, def:physics:phyx-mini-0498:phyxminiproblems-problemphyxmini0498-matchesphotoelectricscenario, def:physics:phyx-mini-0498:phyxminiproblems-problemphyxmini0498-matchessuppliedstoppingpotentialfigure, def:physics:phyx-mini-0498:phyxminiproblems-problemphyxmini0498-hasphysicalphotoelectricparameters, def:physics:phyx-mini-0498:phyxminiproblems-problemphyxmini0498-satisfiesphotoelectriclaws, def:physics:phyx-mini-0498:phyxminiproblems-problemphyxmini0498-usestextbookelementarycharge, def:physics:phyx-mini-0498:phyxminiproblems-problemphyxmini0498-hasadditionaltenmicroamperephotocurrentmeasurement, def:physics:phyx-mini-0498:phyxminiproblems-problemphyxmini0498-displayedphotoelectronratepersecond, def:physics:phyx-mini-0498:phyxminiproblems-problemphyxmini0498-recordeddatasetanswer, def:physics:phyx-mini-0498:phyxminiproblems-problemphyxmini0498-isuniquematchingdisplayedrate}
The assigned autoformalize agent should translate this physics problem into Lean declarations in the covered file, with theorem and lemma proof bodies written as `by sorry`.
\end{theorem}
\begin{proof}
This is an autoformalization task, not a proof task. Produce statements that can later be proved by the physics prover without weakening the physical model.
\end{proof}

% --- Archon declaration topology begin ---
\section{Declaration topology}
\begin{definition}[energy Dimension]
  \label{def:physics:phyx-mini-0498:phyxminiproblems-problemphyxmini0498-energydimension}
  \lean{PhyXMiniProblems.ProblemPhyXMini0498.energyDimension}
  The MLTQ dimension of energy
\end{definition}
\begin{proof}
  This is the declared model definition of energy Dimension.
\end{proof}
\begin{definition}[electric Potential Dimension]
  \label{def:physics:phyx-mini-0498:phyxminiproblems-problemphyxmini0498-electricpotentialdimension}
  \lean{PhyXMiniProblems.ProblemPhyXMini0498.electricPotentialDimension}
  \uses{def:physics:phyx-mini-0498:phyxminiproblems-problemphyxmini0498-energydimension}
  Electric potential has the dimension energy per charge
\end{definition}
\begin{proof}
  This is the declared model definition of electric Potential Dimension.
\end{proof}
\begin{definition}[electric Current Dimension]
  \label{def:physics:phyx-mini-0498:phyxminiproblems-problemphyxmini0498-electriccurrentdimension}
  \lean{PhyXMiniProblems.ProblemPhyXMini0498.electricCurrentDimension}
  Electric current has the dimension charge per time
\end{definition}
\begin{proof}
  This is the declared model definition of electric Current Dimension.
\end{proof}
\begin{definition}[action Dimension]
  \label{def:physics:phyx-mini-0498:phyxminiproblems-problemphyxmini0498-actiondimension}
  \lean{PhyXMiniProblems.ProblemPhyXMini0498.actionDimension}
  \uses{def:physics:phyx-mini-0498:phyxminiproblems-problemphyxmini0498-energydimension}
  Planck's constant has the dimension energy times time
\end{definition}
\begin{proof}
  This is the declared model definition of action Dimension.
\end{proof}
\begin{definition}[Frequency Quantity]
  \label{def:physics:phyx-mini-0498:phyxminiproblems-problemphyxmini0498-frequencyquantity}
  \lean{PhyXMiniProblems.ProblemPhyXMini0498.FrequencyQuantity}
  A nonnegative physical frequency
\end{definition}
\begin{proof}
  This is the declared model definition of Frequency Quantity.
\end{proof}
\begin{definition}[Electric Potential Quantity]
  \label{def:physics:phyx-mini-0498:phyxminiproblems-problemphyxmini0498-electricpotentialquantity}
  \lean{PhyXMiniProblems.ProblemPhyXMini0498.ElectricPotentialQuantity}
  \uses{def:physics:phyx-mini-0498:phyxminiproblems-problemphyxmini0498-electricpotentialdimension}
  A signed physical stopping potential
\end{definition}
\begin{proof}
  This is the declared model definition of Electric Potential Quantity.
\end{proof}
\begin{definition}[Charge Magnitude Quantity]
  \label{def:physics:phyx-mini-0498:phyxminiproblems-problemphyxmini0498-chargemagnitudequantity}
  \lean{PhyXMiniProblems.ProblemPhyXMini0498.ChargeMagnitudeQuantity}
  A nonnegative physical magnitude of electric charge
\end{definition}
\begin{proof}
  This is the declared model definition of Charge Magnitude Quantity.
\end{proof}
\begin{definition}[Electric Current Quantity]
  \label{def:physics:phyx-mini-0498:phyxminiproblems-problemphyxmini0498-electriccurrentquantity}
  \lean{PhyXMiniProblems.ProblemPhyXMini0498.ElectricCurrentQuantity}
  \uses{def:physics:phyx-mini-0498:phyxminiproblems-problemphyxmini0498-electriccurrentdimension}
  A nonnegative physical electric current
\end{definition}
\begin{proof}
  This is the declared model definition of Electric Current Quantity.
\end{proof}
\begin{definition}[Photoelectron Rate Quantity]
  \label{def:physics:phyx-mini-0498:phyxminiproblems-problemphyxmini0498-photoelectronratequantity}
  \lean{PhyXMiniProblems.ProblemPhyXMini0498.PhotoelectronRateQuantity}
  A nonnegative number of emitted photoelectrons per unit time
\end{definition}
\begin{proof}
  This is the declared model definition of Photoelectron Rate Quantity.
\end{proof}
\begin{definition}[Planck Constant Quantity]
  \label{def:physics:phyx-mini-0498:phyxminiproblems-problemphyxmini0498-planckconstantquantity}
  \lean{PhyXMiniProblems.ProblemPhyXMini0498.PlanckConstantQuantity}
  \uses{def:physics:phyx-mini-0498:phyxminiproblems-problemphyxmini0498-actiondimension}
  A nonnegative physical value of Planck's constant
\end{definition}
\begin{proof}
  This is the declared model definition of Planck Constant Quantity.
\end{proof}
\begin{definition}[Work Function Quantity]
  \label{def:physics:phyx-mini-0498:phyxminiproblems-problemphyxmini0498-workfunctionquantity}
  \lean{PhyXMiniProblems.ProblemPhyXMini0498.WorkFunctionQuantity}
  A physical work function, constrained to be nonnegative by the setup premise
\end{definition}
\begin{proof}
  This is the declared model definition of Work Function Quantity.
\end{proof}
\begin{definition}[nonnegative SIReadout]
  \label{def:physics:phyx-mini-0498:phyxminiproblems-problemphyxmini0498-nonnegativesireadout}
  \lean{PhyXMiniProblems.ProblemPhyXMini0498.nonnegativeSIReadout}
  Coherent-SI readout of a nonnegative dimensionful quantity
\end{definition}
\begin{proof}
  This is the declared model definition of nonnegative SIReadout.
\end{proof}
\begin{definition}[signed SIReadout]
  \label{def:physics:phyx-mini-0498:phyxminiproblems-problemphyxmini0498-signedsireadout}
  \lean{PhyXMiniProblems.ProblemPhyXMini0498.signedSIReadout}
  Coherent-SI readout of a signed dimensionful quantity
\end{definition}
\begin{proof}
  This is the declared model definition of signed SIReadout.
\end{proof}
\begin{definition}[frequency In Hertz]
  \label{def:physics:phyx-mini-0498:phyxminiproblems-problemphyxmini0498-frequencyinhertz}
  \lean{PhyXMiniProblems.ProblemPhyXMini0498.frequencyInHertz}
  \uses{def:physics:phyx-mini-0498:phyxminiproblems-problemphyxmini0498-frequencyquantity, def:physics:phyx-mini-0498:phyxminiproblems-problemphyxmini0498-nonnegativesireadout}
  Frequency readout in hertz
\end{definition}
\begin{proof}
  This is the declared model definition of frequency In Hertz.
\end{proof}
\begin{definition}[potential In Volts]
  \label{def:physics:phyx-mini-0498:phyxminiproblems-problemphyxmini0498-potentialinvolts}
  \lean{PhyXMiniProblems.ProblemPhyXMini0498.potentialInVolts}
  \uses{def:physics:phyx-mini-0498:phyxminiproblems-problemphyxmini0498-electricpotentialquantity, def:physics:phyx-mini-0498:phyxminiproblems-problemphyxmini0498-signedsireadout}
  Stopping-potential readout in volts
\end{definition}
\begin{proof}
  This is the declared model definition of potential In Volts.
\end{proof}
\begin{definition}[charge Magnitude In Coulombs]
  \label{def:physics:phyx-mini-0498:phyxminiproblems-problemphyxmini0498-chargemagnitudeincoulombs}
  \lean{PhyXMiniProblems.ProblemPhyXMini0498.chargeMagnitudeInCoulombs}
  \uses{def:physics:phyx-mini-0498:phyxminiproblems-problemphyxmini0498-chargemagnitudequantity, def:physics:phyx-mini-0498:phyxminiproblems-problemphyxmini0498-nonnegativesireadout}
  Charge-magnitude readout in coulombs
\end{definition}
\begin{proof}
  This is the declared model definition of charge Magnitude In Coulombs.
\end{proof}
\begin{definition}[current In Amperes]
  \label{def:physics:phyx-mini-0498:phyxminiproblems-problemphyxmini0498-currentinamperes}
  \lean{PhyXMiniProblems.ProblemPhyXMini0498.currentInAmperes}
  \uses{def:physics:phyx-mini-0498:phyxminiproblems-problemphyxmini0498-electriccurrentquantity, def:physics:phyx-mini-0498:phyxminiproblems-problemphyxmini0498-nonnegativesireadout}
  Electric-current readout in amperes
\end{definition}
\begin{proof}
  This is the declared model definition of current In Amperes.
\end{proof}
\begin{definition}[photoelectron Rate Per Second]
  \label{def:physics:phyx-mini-0498:phyxminiproblems-problemphyxmini0498-photoelectronratepersecond}
  \lean{PhyXMiniProblems.ProblemPhyXMini0498.photoelectronRatePerSecond}
  \uses{def:physics:phyx-mini-0498:phyxminiproblems-problemphyxmini0498-photoelectronratequantity, def:physics:phyx-mini-0498:phyxminiproblems-problemphyxmini0498-nonnegativesireadout}
  Photoelectron-rate readout in electrons per second
\end{definition}
\begin{proof}
  This is the declared model definition of photoelectron Rate Per Second.
\end{proof}
\begin{definition}[planck Constant In Joule Seconds]
  \label{def:physics:phyx-mini-0498:phyxminiproblems-problemphyxmini0498-planckconstantinjouleseconds}
  \lean{PhyXMiniProblems.ProblemPhyXMini0498.planckConstantInJouleSeconds}
  \uses{def:physics:phyx-mini-0498:phyxminiproblems-problemphyxmini0498-planckconstantquantity, def:physics:phyx-mini-0498:phyxminiproblems-problemphyxmini0498-nonnegativesireadout}
  Planck-constant readout in joule-seconds
\end{definition}
\begin{proof}
  This is the declared model definition of planck Constant In Joule Seconds.
\end{proof}
\begin{definition}[work Function In Joules]
  \label{def:physics:phyx-mini-0498:phyxminiproblems-problemphyxmini0498-workfunctioninjoules}
  \lean{PhyXMiniProblems.ProblemPhyXMini0498.workFunctionInJoules}
  \uses{def:physics:phyx-mini-0498:phyxminiproblems-problemphyxmini0498-workfunctionquantity, def:physics:phyx-mini-0498:phyxminiproblems-problemphyxmini0498-signedsireadout}
  Work-function readout in joules
\end{definition}
\begin{proof}
  This is the declared model definition of work Function In Joules.
\end{proof}
\begin{definition}[physlib Elementary Charge In Coulombs]
  \label{def:physics:phyx-mini-0498:phyxminiproblems-problemphyxmini0498-physlibelementarychargeincoulombs}
  \lean{PhyXMiniProblems.ProblemPhyXMini0498.physlibElementaryChargeInCoulombs}
  Physlib represents an elementary charge as a charge *unit*. This scalar is its exact magnitude relative to the Physlib coulomb unit. It is retained to distinguish the exact SI value from the rounded textbook calibration used by the answer choices
\end{definition}
\begin{proof}
  This is the declared model definition of physlib Elementary Charge In Coulombs.
\end{proof}
\begin{definition}[textbook Elementary Charge In Coulombs]
  \label{def:physics:phyx-mini-0498:phyxminiproblems-problemphyxmini0498-textbookelementarychargeincoulombs}
  \lean{PhyXMiniProblems.ProblemPhyXMini0498.textbookElementaryChargeInCoulombs}
  Textbook elementary-charge magnitude used in the multiple-choice arithmetic
\end{definition}
\begin{proof}
  This is the declared model definition of textbook Elementary Charge In Coulombs.
\end{proof}
\begin{definition}[ten Microamperes In Amperes]
  \label{def:physics:phyx-mini-0498:phyxminiproblems-problemphyxmini0498-tenmicroamperesinamperes}
  \lean{PhyXMiniProblems.ProblemPhyXMini0498.tenMicroamperesInAmperes}
  The independent current measurement needed to recover the recorded answer
\end{definition}
\begin{proof}
  This is the declared model definition of ten Microamperes In Amperes.
\end{proof}
\begin{definition}[Figure Axis]
  \label{def:physics:phyx-mini-0498:phyxminiproblems-problemphyxmini0498-figureaxis}
  \lean{PhyXMiniProblems.ProblemPhyXMini0498.FigureAxis}
  The two coordinate axes shown in the supplied raster
\end{definition}
\begin{proof}
  This is the declared model definition of Figure Axis.
\end{proof}
\begin{definition}[Axis Quantity]
  \label{def:physics:phyx-mini-0498:phyxminiproblems-problemphyxmini0498-axisquantity}
  \lean{PhyXMiniProblems.ProblemPhyXMini0498.AxisQuantity}
  Physical role assigned to each axis
\end{definition}
\begin{proof}
  This is the declared model definition of Axis Quantity.
\end{proof}
\begin{definition}[Axis Unit]
  \label{def:physics:phyx-mini-0498:phyxminiproblems-problemphyxmini0498-axisunit}
  \lean{PhyXMiniProblems.ProblemPhyXMini0498.AxisUnit}
  Unit annotation printed beside an axis
\end{definition}
\begin{proof}
  This is the declared model definition of Axis Unit.
\end{proof}
\begin{definition}[Green Line Point]
  \label{def:physics:phyx-mini-0498:phyxminiproblems-problemphyxmini0498-greenlinepoint}
  \lean{PhyXMiniProblems.ProblemPhyXMini0498.GreenLinePoint}
  Three unambiguous grid intersections on the green line
\end{definition}
\begin{proof}
  This is the declared model definition of Green Line Point.
\end{proof}
\begin{definition}[Stopping Potential Frequency Figure]
  \label{def:physics:phyx-mini-0498:phyxminiproblems-problemphyxmini0498-stoppingpotentialfrequencyfigure}
  \lean{PhyXMiniProblems.ProblemPhyXMini0498.StoppingPotentialFrequencyFigure}
  \uses{def:physics:phyx-mini-0498:phyxminiproblems-problemphyxmini0498-frequencyquantity, def:physics:phyx-mini-0498:phyxminiproblems-problemphyxmini0498-electricpotentialquantity, def:physics:phyx-mini-0498:phyxminiproblems-problemphyxmini0498-figureaxis, def:physics:phyx-mini-0498:phyxminiproblems-problemphyxmini0498-axisquantity, def:physics:phyx-mini-0498:phyxminiproblems-problemphyxmini0498-axisunit, def:physics:phyx-mini-0498:phyxminiproblems-problemphyxmini0498-greenlinepoint}
  Physical quantities represented by the plotted green-line points, together with the visible axis metadata. The Boolean fields record the absence of the two readouts that would be needed for the question as written
\end{definition}
\begin{proof}
  This is the declared model definition of Stopping Potential Frequency Figure.
\end{proof}
\begin{definition}[Quantum Process]
  \label{def:physics:phyx-mini-0498:phyxminiproblems-problemphyxmini0498-quantumprocess}
  \lean{PhyXMiniProblems.ProblemPhyXMini0498.QuantumProcess}
  The physical process represented by the setup
\end{definition}
\begin{proof}
  This is the declared model definition of Quantum Process.
\end{proof}
\begin{definition}[Emitted Particle]
  \label{def:physics:phyx-mini-0498:phyxminiproblems-problemphyxmini0498-emittedparticle}
  \lean{PhyXMiniProblems.ProblemPhyXMini0498.EmittedParticle}
  Species emitted by the illuminated surface
\end{definition}
\begin{proof}
  This is the declared model definition of Emitted Particle.
\end{proof}
\begin{definition}[Photoelectric Experiment]
  \label{def:physics:phyx-mini-0498:phyxminiproblems-problemphyxmini0498-photoelectricexperiment}
  \lean{PhyXMiniProblems.ProblemPhyXMini0498.PhotoelectricExperiment}
  \uses{def:physics:phyx-mini-0498:phyxminiproblems-problemphyxmini0498-frequencyquantity, def:physics:phyx-mini-0498:phyxminiproblems-problemphyxmini0498-electricpotentialquantity, def:physics:phyx-mini-0498:phyxminiproblems-problemphyxmini0498-chargemagnitudequantity, def:physics:phyx-mini-0498:phyxminiproblems-problemphyxmini0498-electriccurrentquantity, def:physics:phyx-mini-0498:phyxminiproblems-problemphyxmini0498-photoelectronratequantity, def:physics:phyx-mini-0498:phyxminiproblems-problemphyxmini0498-planckconstantquantity, def:physics:phyx-mini-0498:phyxminiproblems-problemphyxmini0498-workfunctionquantity, def:physics:phyx-mini-0498:phyxminiproblems-problemphyxmini0498-stoppingpotentialfrequencyfigure, def:physics:phyx-mini-0498:phyxminiproblems-problemphyxmini0498-quantumprocess, def:physics:phyx-mini-0498:phyxminiproblems-problemphyxmini0498-emittedparticle}
  The experiment stores energy-side observables, charge/count-side observables, and the supplied graph. Neither the requested scalar rate nor an answer choice is a field
\end{definition}
\begin{proof}
  This is the declared model definition of Photoelectric Experiment.
\end{proof}
\begin{definition}[Matches Photoelectric Scenario]
  \label{def:physics:phyx-mini-0498:phyxminiproblems-problemphyxmini0498-matchesphotoelectricscenario}
  \lean{PhyXMiniProblems.ProblemPhyXMini0498.MatchesPhotoelectricScenario}
  \uses{def:physics:phyx-mini-0498:phyxminiproblems-problemphyxmini0498-photoelectricexperiment}
  The prose identifies the phenomenon and the emitted particles
\end{definition}
\begin{proof}
  This is the declared model definition of Matches Photoelectric Scenario.
\end{proof}
\begin{definition}[Matches Supplied Stopping Potential Figure]
  \label{def:physics:phyx-mini-0498:phyxminiproblems-problemphyxmini0498-matchessuppliedstoppingpotentialfigure}
  \lean{PhyXMiniProblems.ProblemPhyXMini0498.MatchesSuppliedStoppingPotentialFigure}
  \uses{def:physics:phyx-mini-0498:phyxminiproblems-problemphyxmini0498-frequencyinhertz, def:physics:phyx-mini-0498:phyxminiproblems-problemphyxmini0498-potentialinvolts, def:physics:phyx-mini-0498:phyxminiproblems-problemphyxmini0498-photoelectricexperiment}
  Evidence read directly from image 498. In particular, the line begins at frequency coordinate 1, not at the origin as claimed by the auxiliary caption. No current or count-rate datum occurs in this structure
\end{definition}
\begin{proof}
  This is the declared model definition of Matches Supplied Stopping Potential Figure.
\end{proof}
\begin{definition}[Has Physical Photoelectric Parameters]
  \label{def:physics:phyx-mini-0498:phyxminiproblems-problemphyxmini0498-hasphysicalphotoelectricparameters}
  \lean{PhyXMiniProblems.ProblemPhyXMini0498.HasPhysicalPhotoelectricParameters}
  \uses{def:physics:phyx-mini-0498:phyxminiproblems-problemphyxmini0498-frequencyinhertz, def:physics:phyx-mini-0498:phyxminiproblems-problemphyxmini0498-chargemagnitudeincoulombs, def:physics:phyx-mini-0498:phyxminiproblems-problemphyxmini0498-planckconstantinjouleseconds, def:physics:phyx-mini-0498:phyxminiproblems-problemphyxmini0498-workfunctioninjoules, def:physics:phyx-mini-0498:phyxminiproblems-problemphyxmini0498-photoelectricexperiment}
  Positivity conditions for the physical constants and operating light
\end{definition}
\begin{proof}
  This is the declared model definition of Has Physical Photoelectric Parameters.
\end{proof}
\begin{definition}[Satisfies Photoelectric Laws]
  \label{def:physics:phyx-mini-0498:phyxminiproblems-problemphyxmini0498-satisfiesphotoelectriclaws}
  \lean{PhyXMiniProblems.ProblemPhyXMini0498.SatisfiesPhotoelectricLaws}
  \uses{def:physics:phyx-mini-0498:phyxminiproblems-problemphyxmini0498-frequencyinhertz, def:physics:phyx-mini-0498:phyxminiproblems-problemphyxmini0498-potentialinvolts, def:physics:phyx-mini-0498:phyxminiproblems-problemphyxmini0498-chargemagnitudeincoulombs, def:physics:phyx-mini-0498:phyxminiproblems-problemphyxmini0498-currentinamperes, def:physics:phyx-mini-0498:phyxminiproblems-problemphyxmini0498-photoelectronratepersecond, def:physics:phyx-mini-0498:phyxminiproblems-problemphyxmini0498-planckconstantinjouleseconds, def:physics:phyx-mini-0498:phyxminiproblems-problemphyxmini0498-workfunctioninjoules, def:physics:phyx-mini-0498:phyxminiproblems-problemphyxmini0498-photoelectricexperiment}
  The two governing laws needed to distinguish the graph's energy information from rate information: Einstein's stopping-potential relation e V\_stop = h f - φ, where a nonnegative stopping potential is defined; collected conventional photocurrent equals the elementary charge magnitude times the emitted-electron count rate. Both laws are uniform relations and contain no numerical answer choice
\end{definition}
\begin{proof}
  This is the declared model definition of Satisfies Photoelectric Laws.
\end{proof}
\begin{definition}[Uses Textbook Elementary Charge]
  \label{def:physics:phyx-mini-0498:phyxminiproblems-problemphyxmini0498-usestextbookelementarycharge}
  \lean{PhyXMiniProblems.ProblemPhyXMini0498.UsesTextbookElementaryCharge}
  \uses{def:physics:phyx-mini-0498:phyxminiproblems-problemphyxmini0498-chargemagnitudeincoulombs, def:physics:phyx-mini-0498:phyxminiproblems-problemphyxmini0498-textbookelementarychargeincoulombs, def:physics:phyx-mini-0498:phyxminiproblems-problemphyxmini0498-photoelectricexperiment}
  Standard rounded elementary-charge calibration used by the options. This is a constant calibration, not the requested rate
\end{definition}
\begin{proof}
  This is the declared model definition of Uses Textbook Elementary Charge.
\end{proof}
\begin{definition}[Has Additional Ten Microampere Photocurrent Measurement]
  \label{def:physics:phyx-mini-0498:phyxminiproblems-problemphyxmini0498-hasadditionaltenmicroamperephotocurrentmeasurement}
  \lean{PhyXMiniProblems.ProblemPhyXMini0498.HasAdditionalTenMicroamperePhotocurrentMeasurement}
  \uses{def:physics:phyx-mini-0498:phyxminiproblems-problemphyxmini0498-currentinamperes, def:physics:phyx-mini-0498:phyxminiproblems-problemphyxmini0498-tenmicroamperesinamperes, def:physics:phyx-mini-0498:phyxminiproblems-problemphyxmini0498-photoelectricexperiment}
  An additional photocurrent measurement which is required for the recorded answer but is absent from image 498. Naming it separately prevents this datum from being mistaken for a figure readout
\end{definition}
\begin{proof}
  This is the declared model definition of Has Additional Ten Microampere Photocurrent Measurement.
\end{proof}
\begin{definition}[Answer Choice]
  \label{def:physics:phyx-mini-0498:phyxminiproblems-problemphyxmini0498-answerchoice}
  \lean{PhyXMiniProblems.ProblemPhyXMini0498.AnswerChoice}
  Labels printed beside the four answer choices
\end{definition}
\begin{proof}
  This is the declared model definition of Answer Choice.
\end{proof}
\begin{definition}[displayed Photoelectron Rate Per Second]
  \label{def:physics:phyx-mini-0498:phyxminiproblems-problemphyxmini0498-displayedphotoelectronratepersecond}
  \lean{PhyXMiniProblems.ProblemPhyXMini0498.displayedPhotoelectronRatePerSecond}
  \uses{def:physics:phyx-mini-0498:phyxminiproblems-problemphyxmini0498-answerchoice}
  Photoelectron-rate value printed beside each answer label
\end{definition}
\begin{proof}
  This is the declared model definition of displayed Photoelectron Rate Per Second.
\end{proof}
\begin{definition}[recorded Dataset Answer]
  \label{def:physics:phyx-mini-0498:phyxminiproblems-problemphyxmini0498-recordeddatasetanswer}
  \lean{PhyXMiniProblems.ProblemPhyXMini0498.recordedDatasetAnswer}
  \uses{def:physics:phyx-mini-0498:phyxminiproblems-problemphyxmini0498-answerchoice}
  The source dataset's recorded answer, retained only as metadata
\end{definition}
\begin{proof}
  This is the declared model definition of recorded Dataset Answer.
\end{proof}
\begin{definition}[Is Unique Matching Displayed Rate]
  \label{def:physics:phyx-mini-0498:phyxminiproblems-problemphyxmini0498-isuniquematchingdisplayedrate}
  \lean{PhyXMiniProblems.ProblemPhyXMini0498.IsUniqueMatchingDisplayedRate}
  \uses{def:physics:phyx-mini-0498:phyxminiproblems-problemphyxmini0498-photoelectronratepersecond, def:physics:phyx-mini-0498:phyxminiproblems-problemphyxmini0498-photoelectricexperiment, def:physics:phyx-mini-0498:phyxminiproblems-problemphyxmini0498-answerchoice, def:physics:phyx-mini-0498:phyxminiproblems-problemphyxmini0498-displayedphotoelectronratepersecond}
  A choice is the unique printed value equal to the experiment's rate
\end{definition}
\begin{proof}
  This is the declared model definition of Is Unique Matching Displayed Rate.
\end{proof}
\begin{lemma}[supplied Figure has no current or rate readout]
  \label{lem:physics:phyx-mini-0498:phyxminiproblems-problemphyxmini0498-suppliedfigure-has-no-current-or-rate-readout}
  \lean{PhyXMiniProblems.ProblemPhyXMini0498.suppliedFigure_has_no_current_or_rate_readout}
  \uses{def:physics:phyx-mini-0498:phyxminiproblems-problemphyxmini0498-photoelectricexperiment, def:physics:phyx-mini-0498:phyxminiproblems-problemphyxmini0498-matchessuppliedstoppingpotentialfigure}
  The supplied stopping-potential graph explicitly contains no rate datum
\end{lemma}
\begin{proof}
  The conclusion follows from the governing relations and figure readouts named in the statement of supplied Figure has no current or rate readout.
\end{proof}
\begin{lemma}[photoelectron Rate of ten Microampere Measurement]
  \label{lem:physics:phyx-mini-0498:phyxminiproblems-problemphyxmini0498-photoelectronrate-of-tenmicroamperemeasurement}
  \lean{PhyXMiniProblems.ProblemPhyXMini0498.photoelectronRate_of_tenMicroampereMeasurement}
  \uses{def:physics:phyx-mini-0498:phyxminiproblems-problemphyxmini0498-photoelectronratepersecond, def:physics:phyx-mini-0498:phyxminiproblems-problemphyxmini0498-photoelectricexperiment, def:physics:phyx-mini-0498:phyxminiproblems-problemphyxmini0498-hasphysicalphotoelectricparameters, def:physics:phyx-mini-0498:phyxminiproblems-problemphyxmini0498-satisfiesphotoelectriclaws, def:physics:phyx-mini-0498:phyxminiproblems-problemphyxmini0498-usestextbookelementarycharge, def:physics:phyx-mini-0498:phyxminiproblems-problemphyxmini0498-hasadditionaltenmicroamperephotocurrentmeasurement}
  With the missing current measurement made explicit, the charge-rate law and the rounded elementary charge give 6.25 × 10¹³ electrons per second
\end{lemma}
\begin{proof}
  The conclusion follows from the governing relations and figure readouts named in the statement of photoelectron Rate of ten Microampere Measurement.
\end{proof}
\begin{definition}[Uses the Physlib elementary charge]
  \label{def:physics:phyx-mini-0498:phyxminiproblems-problemphyxmini0498-usesphyslibelementarycharge}
  \lean{PhyXMiniProblems.ProblemPhyXMini0498.UsesPhyslibElementaryCharge}
  \uses{def:physics:phyx-mini-0498:phyxminiproblems-problemphyxmini0498-photoelectricexperiment, def:physics:phyx-mini-0498:phyxminiproblems-problemphyxmini0498-chargemagnitudeincoulombs, def:physics:phyx-mini-0498:phyxminiproblems-problemphyxmini0498-physlibelementarychargeincoulombs}
  The experiment's electron-charge magnitude is calibrated to Physlib's exact elementary-charge magnitude, without assuming any current or count rate.
\end{definition}
\begin{proof}
  This is the exact physical-constant calibration.
\end{proof}
\begin{definition}[Compatible current and rate readouts]
  \label{def:physics:phyx-mini-0498:phyxminiproblems-problemphyxmini0498-compatiblecurrentratereadouts}
  \lean{PhyXMiniProblems.ProblemPhyXMini0498.CompatibleCurrentRateReadouts}
  \uses{def:physics:phyx-mini-0498:phyxminiproblems-problemphyxmini0498-photoelectricexperiment, def:physics:phyx-mini-0498:phyxminiproblems-problemphyxmini0498-chargemagnitudeincoulombs}
  A hypothetical nonnegative current-rate pair is compatible when it satisfies the collected-charge law $I=eN$.
\end{definition}
\begin{proof}
  This is the governing compatibility relation, not an assertion that either readout was measured.
\end{proof}
\begin{definition}[Displayed rate has a compatible unmeasured current]
  \label{def:physics:phyx-mini-0498:phyxminiproblems-problemphyxmini0498-displayedratehascompatibleunmeasuredcurrent}
  \lean{PhyXMiniProblems.ProblemPhyXMini0498.DisplayedRateHasCompatibleUnmeasuredCurrent}
  \uses{def:physics:phyx-mini-0498:phyxminiproblems-problemphyxmini0498-photoelectricexperiment, def:physics:phyx-mini-0498:phyxminiproblems-problemphyxmini0498-answerchoice, def:physics:phyx-mini-0498:phyxminiproblems-problemphyxmini0498-displayedphotoelectronratepersecond, def:physics:phyx-mini-0498:phyxminiproblems-problemphyxmini0498-compatiblecurrentratereadouts}
  A displayed rate is compatible when some nonnegative, unmeasured current can accompany it under $I=eN$.
\end{definition}
\begin{proof}
  This is the existential compatibility definition and does not claim the rate is actual.
\end{proof}
\begin{lemma}[Every displayed rate has a compatible unmeasured current]
  \label{lem:physics:phyx-mini-0498:phyxminiproblems-problemphyxmini0498-everydisplayedrate-has-compatible-unmeasuredcurrent}
  \lean{PhyXMiniProblems.ProblemPhyXMini0498.everyDisplayedRate_has_compatible_unmeasuredCurrent}
  \uses{def:physics:phyx-mini-0498:phyxminiproblems-problemphyxmini0498-photoelectricexperiment, def:physics:phyx-mini-0498:phyxminiproblems-problemphyxmini0498-hasphysicalphotoelectricparameters, def:physics:phyx-mini-0498:phyxminiproblems-problemphyxmini0498-usesphyslibelementarycharge, def:physics:phyx-mini-0498:phyxminiproblems-problemphyxmini0498-answerchoice, def:physics:phyx-mini-0498:phyxminiproblems-problemphyxmini0498-displayedratehascompatibleunmeasuredcurrent}
  Each printed rate is compatible with some possible unmeasured photocurrent, so the supplied stopping-potential graph cannot select a numeric rate.
\end{lemma}
\begin{proof}
  For an arbitrary displayed rate $N$, choose the candidate current $I=eN$.  Positivity of the exact charge and nonnegativity of every printed rate give the two sign conditions, and the charge-rate equation holds by construction.
\end{proof}
% --- Archon declaration topology end ---
% --- Archon physics formalization source end ---
